\begin{definition}[Finite-field character and Gauss-sum API]\label{mod:finitefield-characterapi}
For a field $k$, put
\[
 \mathrm{FiniteAddChar}(k)=\operatorname{AddChar}(k,\mathbf C),
 \qquad
 \mathrm{FiniteMulChar}(k)=\operatorname{MulChar}(k,\mathbf C).
\]
Here Mathlib's multiplicative character is the canonical extension by zero
of a homomorphism $k^\times\to\mathbf C^\times$.

Suppose now that $k$ is finite.  If $\psi\neq1$ and $\chi\neq1$, the
exported orthogonality statements are
\[
 \sum_{x\in k}\psi(x)=0,
 \qquad
 \sum_{x\in k}\chi(x)=0.
\]
Mathlib's ordinary Gauss sum is identified with the manuscript's
$k^\times$-indexed sum:
\[
 G_k(\chi,\psi)=\operatorname{gaussSum}(\chi,\psi)
 =\sum_{x\in k^\times}\chi(x)\psi(x).
\]
Langlands's normalized sum is defined, with the manuscript's sign and
inverse convention, by
\[
 \operatorname{langlandsGaussSum}(\chi,\psi)
 =\tau_k(\chi,\psi)
 :=-G_k(\chi^{-1},\psi)
 =-\sum_{x\in k^\times}\chi(x)^{-1}\psi(x).
\]
For $\psi\neq1$ the API proves
\[
 G_k(1,\psi)=-1,
 \qquad \tau_k(1,\psi)=1,
 \qquad G_k(\chi,\psi)\neq0,
 \qquad \tau_k(\chi,\psi)\neq0
\]
for every $\chi$.  If also $\chi\neq1$, it proves both the conjugate
product and real norm forms of the precise magnitude identity:
\[
 G_k(\chi,\psi)\overline{G_k(\chi,\psi)}=|k|,
 \qquad \lVert G_k(\chi,\psi)\rVert^2=|k|.
\]
For $a\in k^\times$ and $\psi_a(x)=\psi(ax)$, the scaling declarations
are
\[
 G_k(\chi,\psi_a)=\chi(a)^{-1}G_k(\chi,\psi),
 \qquad
 \tau_k(\chi,\psi_a)=\chi(a)\tau_k(\chi,\psi).
\]

Finally, for a finite field extension $K/k$, the definitions
$\operatorname{traceAddChar}_{K/k}(\psi)=\psi\circ\operatorname{Tr}_{K/k}$
and
$\operatorname{normMulChar}_{K/k}(\chi)=\chi\circ N_{K/k}$ are accompanied
by their pointwise evaluation lemmas and by
\[
 \operatorname{traceAddChar}_{K/k}(\psi)\neq1\iff\psi\neq1,
 \qquad
 \operatorname{normMulChar}_{K/k}(\chi)\neq1\iff\chi\neq1,
\]
together with
$\operatorname{normMulChar}_{K/k}(\chi^{-1})
 =\operatorname{normMulChar}_{K/k}(\chi)^{-1}$.
\LeanFile{LanglandsFirstMainLemma/FiniteField/CharacterAPI.lean}
\lean{LanglandsFirstMainLemma.langlandsGaussSum}
\leanok
\SourceText{Lemma lem:finite-gauss-magnitude and Theorem thm:HD-lift.}
\uses{mod:basic-finiteproducts,mod:localfield-residuefield}
\FormalizationContract
\end{definition}
